\scanpage{96}
\[ F^{\mathrm{opp}}:B\longrightarrow(\mathrm{Cat}) \]

défini par

\[
\begin{aligned}
F^{\mathrm{opp}}(x)&=F(x)^\circ\\
F^{\mathrm{opp}}(f^{\unsure{\circ|'|{}}})&=F(f)^\circ
\end{aligned}
\]

\deleted{\unsure{et|lui}}

et formule \unsure{analogue} pour isom de transitivité.

Il n'est pas évident qu'il en existe une description intrinsèque à \uline{isom} près, en évitant \unsure{tt|tout} recours au pseudo-foncteur.

Pour que $(\deleted{\unsure{B}}{\unsure{E|\mathcal E}}/B)^{\mathrm{opp}}$ soit aussi cofibrée, il f. et s. que les foncteurs $f^*:E_y\mathrel{\rule{1.4em}{.35pt}}E_x$ aient des adjoints à \uline{droite} (condition opposée à celle pour que $E/B$ soit cofibrée).

\uline{À retenir} le $=$ pseudo-foncteur définit deux catégories (fibrées) bien distinctes sur $B$.

% archive sheet number: 95
